% =============================================================
% 开题报告正文
% 论文题目：基于物流领域大语言模型的智能运营分析助手设计与实现
% 作者：巩德志  学号：22460420  导师：张朋
% =============================================================

% -------------------------------------------------------------
% 第一章 选题背景及意义
% -------------------------------------------------------------
\chapter{选题背景及意义}

\section{研究背景}

% TODO: 物流运营场景痛点、大模型应用趋势

\section{研究意义}

% TODO: 业务价值 + 学术价值

\section{本文主要工作}

% TODO: 简述两个算法创新 + 系统设计与实现


% -------------------------------------------------------------
% 第二章 国内外研究现状
% -------------------------------------------------------------
\chapter{国内外研究现状}

\section{领域大语言模型研究现状}

% TODO

\section{面向智能体的工具使用与规划研究现状}

% TODO

\section{检索增强生成研究现状}

% TODO

\section{代价敏感学习与安全拒答研究现状}

% TODO

\section{小结与不足}

% TODO


% -------------------------------------------------------------
% 第三章 主要研究方法
% -------------------------------------------------------------
\chapter{主要研究方法}

\section{大语言模型微调方法}

% TODO: SFT / LoRA / QLoRA / DPO

\section{面向智能体的工具编排方法}

% TODO: ReAct / Function Calling / FSM 编排

\section{检索增强生成方法}

% TODO: RAG / Dense Retrieval / Rerank

\section{代价敏感学习与不确定性估计方法}

% TODO: 代价矩阵 / Focal Loss / 温度标定 / 双阈值拒答

\section{Agent 系统设计方法}

% TODO: 分层架构 / 记忆 / 工具抽象 / 会话状态


% -------------------------------------------------------------
% 第四章 研究内容与核心创新点
% -------------------------------------------------------------
\chapter{研究内容与核心创新点}

\section{代价敏感的分层路由与安全拒答机制}

\subsection{问题定义}

% TODO: 成本与风险失衡问题

\subsection{方法设计}

% TODO: RoBERTa-base 分类头 + 代价矩阵加权 Focal Loss + 温度标定 + 双阈值拒答

\subsection{关键技术点}

% TODO: 代价矩阵构造 / LoRA 增量学习 / 置信度校准

\subsection{评测方案}

% TODO: 路由准确率 / 拒答率 / 期望代价 / 端到端任务成功率


\section{不确定性感知的工具规划与自修复框架}

\subsection{问题定义}

% TODO: 工具调用幻觉、参数乱填、失败死循环

\subsection{方法设计}

% TODO: Qwen2.5-7B QLoRA SFT+DPO 规划器 + 1.5B verifier + FSM 修复

\subsection{关键技术点}

% TODO: 偏好数据构造 / verifier 训练目标 / FSM 修复策略

\subsection{评测方案}

% TODO: 工具调用准确率 / 参数正确率 / 自修复成功率 / 任务完成率


% -------------------------------------------------------------
% 第五章 系统设计与实现
% -------------------------------------------------------------
\chapter{系统设计与实现}

\section{系统总体架构}

% TODO: 四层架构 (会话层 / 规划层 / 工具层 / 知识层),模块边界与交互

\section{多轮对话与会话状态管理}

% TODO: 上下文组织 / 长期记忆 / 追问澄清 / 意图延续

\section{规划器与执行器设计}

% TODO: LLM 规划调用流程 / 执行轨迹 / 依赖调度 / 失败处理

\section{工具系统设计与实现}

% TODO: 工具抽象 / 参数 schema / 调用协议 / 结果回填

\section{运营场景工具集实现}

% TODO: 指标查询 / 异常定位 / SOP 检索 / 报表生成

\section{知识检索子系统实现}

% TODO: 领域知识库组织 / 向量+关键词融合检索 / 结果排序

\section{分层路由与拒答模块的系统集成}

% TODO: 算法模型在 Agent 链路中的落地

\section{业务侧应用集成}

% TODO: 会话界面 / 权限控制 / 反馈回收


% -------------------------------------------------------------
% 第六章 研究方案与进度安排
% -------------------------------------------------------------
\chapter{研究方案与进度安排}

\section{技术路线}

% TODO

\section{实验与数据准备}

% TODO

\section{进度安排}

% TODO: 甘特图

\section{可能的困难与对策}

% TODO


% -------------------------------------------------------------
% 第七章 预期成果
% -------------------------------------------------------------
\chapter{预期成果}

\section{系统成果}

% TODO

\section{学术成果}

% TODO: 论文 / 专利

\section{应用价值}

% TODO
